% Лабораторная работа № 6. Компиляция: make (или pdflatex report_lab6.tex дважды).
\documentclass[a4paper,14pt]{extarticle}

\usepackage[utf8]{inputenc}
\usepackage[T2A]{fontenc}
\usepackage[russian]{babel}
\usepackage{geometry}
\usepackage{setspace}
\usepackage{indentfirst}
\usepackage{titlesec}
\usepackage{tocloft}
\usepackage{enumitem}
\usepackage{booktabs}
\usepackage{array}
\usepackage{longtable}
\usepackage{graphicx}
\usepackage{fancyvrb}
\usepackage[hidelinks,unicode]{hyperref}

\geometry{left=30mm,right=15mm,top=20mm,bottom=20mm}
\onehalfspacing
\setlength{\parindent}{12.5mm}
\setlength{\parskip}{0pt}
\renewcommand{\contentsname}{СОДЕРЖАНИЕ}
\renewcommand{\refname}{СПИСОК ИСПОЛЬЗОВАННЫХ ИСТОЧНИКОВ}
\titleformat{\section}{\centering\bfseries\fontsize{14pt}{16pt}\selectfont}{\thesection}{1em}{}
\titleformat{\subsection}{\bfseries\fontsize{14pt}{16pt}\selectfont}{\thesubsection}{1em}{}
\titlespacing*{\section}{0pt}{12pt}{12pt}
\titlespacing*{\subsection}{\parindent}{12pt}{6pt}
\setlist{nosep}
\setcounter{tocdepth}{2}
\hypersetup{pdftitle={Отчёт по лабораторной работе № 6},pdfauthor={Студент группы P3113}}

\begin{document}
\begin{titlepage}
\begin{center}
Федеральное государственное автономное образовательное учреждение\\
высшего образования\\
\textbf{«Национальный исследовательский университет ИТМО»}\\[10mm]

Факультет программной инженерии и компьютерной техники\\
Дисциплина «Программирование»\\[34mm]

{\large Лабораторная работа № 6}\\[4mm]
Клиент-серверное приложение\\
Вариант \textbf{9585}
\end{center}

\vfill
\begin{flushright}
Выполнил:\\
Студент группы P3113
Ушаков Роман Павлович\\[5mm]
Проверил:\\
Гиря Максим Дмитриевич

\end{flushright}

\vfill
\begin{center}
г. Санкт-Петербург\\ 2026
\end{center}
\end{titlepage}

\setcounter{page}{2}
\tableofcontents
\newpage

\section{Задание}

Разделить программу из лабораторной работы №5 на клиентский и серверный модули. Серверный модуль должен осуществлять выполнение команд по управлению коллекцией. Клиентский модуль должен в интерактивном режиме считывать команды, передавать их для выполнения на сервер и выводить результаты выполнения.


\section{Цель работы}

Освоить построение простого клиент-серверного приложения на Java, передачу сериализованных объектов по UDP, использование неблокирующего ввода-вывода и разделение ответственности между клиентом, сервером и общим модулем.

\section{Используемые средства}

Язык реализации --- Java~17. Сборка выполняется Apache Maven. Для сохранения коллекции применяется библиотека Gson. В качестве транспорта используются классы \texttt{DatagramChannel} и \texttt{Selector} из пакета \texttt{java.nio.channels}.

\section{Диаграмма классов UML}

UML-диаграмма основных классов вынесена в отдельный файл \href{uml_class_diagram.pdf}{\texttt{uml\_class\_diagram.pdf}}. Исходник диаграммы в PlantUML находится в файле \texttt{uml\_class\_diagram.puml} рядом с данным документом.

\section{Исходный код программы}

Ввиду немалого размера проекта, его код и более подробное описание можно посмотреть на \href{https://gitlab.se.ifmo.ru/amoruch/programming_lab6}{gitlab}.

\section{Результат работы программы}


\end{document}
